\chapter{Metody MCMC: łańcuchy Markowa, estymacja i diagnostyka}
\label{ch:mcmc}
\paragraph{Powiązania i odsyłacze.} Podstawy łańcuchów Markowa omawiają tom I, rozdział~\ref{I-ch:markov}, i tom II, rozdział~\ref{II-ch:ctmc}; metodę Monte Carlo — rozdział~\ref{ch:monte-carlo}.


\section{Gdy nie umiemy losować z rozkładu a posteriori}
Niech dane $x$ wyznaczają posterior o gęstości
$\pi(\theta\mid x)=\widetilde\pi(\theta\mid x)/Z(x)$,
gdzie $\widetilde\pi\ge0$ można obliczać punktowo, lecz
$Z(x)=\int_\Theta\widetilde\pi(\theta\mid x)\,
\dd\lambda_d(\theta)$ nie ma wygodnego wzoru.
Zakładamy $0<Z(x)<\infty$, ponieważ niewłaściwy posterior
nie jest rozkładem, z którego można poprawnie losować.
Metody \emph{Markov chain Monte Carlo} (MCMC) konstruują
łańcuch o rozkładzie niezmienniczym $\pi$ i estymują
$\mu=\int_\Theta f(\theta)\pi(\theta\mid x)\,
\dd\lambda_d(\theta)$ przez średnią wzdłuż trajektorii.
Nie należy zakładać, że kolejne stany łańcucha są niezależne.
Związek długookresowej wariancji ze zjawiskiem autokorelacji omówiliśmy wcześniej w rozdziale~\ref{ch:time-series}.
Wszystkie te pojęcia odnoszą się do ustalonych danych $x$;
parametr $N$ poniżej oznacza długość łańcucha, nie liczbę
danych wykorzystanych do obliczenia posterioru.

\begin{definicja}[Jądro Markowa i miara niezmiennicza]
\label{def:mcmc-kernel}
Na przestrzeni mierzalnej $(\Theta,\mathcal A)$ \emph{jądrem Markowa}
nazywamy odwzorowanie $P\colon\Theta\times\mathcal A\to[0,1]$
takie, że $P(\theta,\cdot)$ jest prawdopodobieństwem dla
każdego $\theta$ oraz $P(\cdot,A)$ jest mierzalna dla każdego $A$.
Miara probabilistyczna $\pi$ jest \emph{niezmiennicza}, jeśli
$\int_\Theta P(\theta,A)\,\dd\pi(\theta)=\pi(A)$ dla każdego $A$.
Jądro spełnia \emph{warunek szczegółowej równowagi} względem $\pi$,
gdy miary na $\Theta\times\Theta$ zadane przez
$\pi(\dd\theta)P(\theta,\dd\vartheta)$ i
$\pi(\dd\vartheta)P(\vartheta,\dd\theta)$ są równe.
Taki łańcuch nazywamy \emph{odwracalnym} względem $\pi$.
\end{definicja}

\begin{stwierdzenie}[Szczegółowa równowaga implikuje niezmienniczość]
\label{prop:balance-invariant}
Jeżeli jądro $P$ spełnia szczegółową równowagę względem
prawdopodobieństwa $\pi$, to $\pi P=\pi$.
\end{stwierdzenie}
\begin{proof}
Dla każdego $A\in\mathcal A$ całkujemy równość miar
szczegółowej równowagi po $\Theta\times A$.
Lewa strona jest $\int_\Theta P(\theta,A)\,\dd\pi(\theta)$,
a prawa $\int_A P(\vartheta,\Theta)\,\dd\pi(\vartheta)=\pi(A)$.
\end{proof}

\begin{uwaga}[Niezmienniczość nie jest jeszcze zbieżnością]
Łańcuch $X_{k+1}=X_k$ ma \emph{każdy} rozkład jako
niezmienniczy, ale nie bada rozkładu docelowego.
Aby średnie MCMC były użyteczne, potrzebujemy dodatkowo
odpowiednich własności ergodycznych: przede wszystkim
osiągalności regionów o dodatniej masie docelowej,
a w typowych twierdzeniach również nieokresowości.
\end{uwaga}

\section{Metropolis--Hastings: dowód poprawności algorytmu}
\begin{twierdzenie}[Jądro Metropolisa--Hastingsa]
\label{thm:mh-detailed}
Niech $\Theta\subseteq\R^d$ będzie mierzalna,
$\pi$ gęstością prawdopodobieństwa względem $\lambda_d$,
a $q(y\mid x)$ mierzalną gęstością propozycji względem
$\lambda_d$ dla każdego $x\in\Theta$.
Dla $\pi(x)q(y\mid x)>0$ zdefiniujmy
\[
 \alpha(x,y)=\min\left\{1,
 \frac{\pi(y)q(x\mid y)}{\pi(x)q(y\mid x)}\right\};
\]
na pozostałych parach z $q(y\mid x)>0$ można ustalić
$\alpha(x,y)$ w dowolny mierzalny sposób zgodny z
zachowaniem na nośniku $\pi$; w szczególności dla
$\pi(x)>0$ i $\pi(y)=0$ przyjmujemy $\alpha(x,y)=0$.
Algorytm: z punktu $x$ losujemy $Y\sim q(\cdot\mid x)$,
a następnie przechodzimy do $Y$ z prawdopodobieństwem
$\alpha(x,Y)$, w przeciwnym razie pozostajemy w $x$.
Tak otrzymane jądro $P$ ma $\pi$ jako rozkład niezmienniczy.
Do obliczania $\alpha$ wystarczy nieunormowana
gęstość $\widetilde\pi=Z\pi$, gdy $0<Z<\infty$.
\end{twierdzenie}
\begin{proof}
Dla $x\ne y$ gęstość przepływu zaakceptowanych przejść
względem $\lambda_{2d}$ jest równa
\[
 \pi(x)q(y\mid x)\alpha(x,y)
 =\min\{\pi(x)q(y\mid x),\pi(y)q(x\mid y)\},
\]
co jest symetryczne w $x,y$. Pozostała masa przejść
znajduje się na przekątnej $\{(x,x):x\in\Theta\}$,
a ta miara także pozostaje niezmieniona po zamianie
współrzędnych. Dla stanów spoza nośnika $\pi$
wybór jądra nie wpływa na powyższe równości miar.
Otrzymaliśmy szczegółową równowagę i możemy zastosować
stwierdzenie~\ref{prop:balance-invariant}.
Stała $Z$ skraca się w ilorazie $\pi(y)/\pi(x)$.
\end{proof}

\begin{przyklad}[Dwie propozycje i ich konsekwencje]
Jeżeli $q(y\mid x)=q(x\mid y)$, to
$\alpha(x,y)=\min\{1,\widetilde\pi(y)/\widetilde\pi(x)\}$.
Tak jest np. dla propozycji $Y=X+\varepsilon$,
$\varepsilon\sim\Normal_d(0,s^2I_d)$; metoda nazywa się
\emph{random-walk Metropolis}. Dla małego $s$ akceptacja
może być wysoka, lecz łańcuch porusza się powoli. Dla
bardzo dużego $s$ częste odrzucenia mogą powodować
powtarzanie tych samych stanów.
Jeżeli natomiast $q(y\mid x)=g(y)$ nie zależy od $x$,
to $\alpha(x,y)=\min\{1,\pi(y)g(x)/(\pi(x)g(y))\}$.
Nie wolno pominąć ilorazu propozycji, chyba że
propozycja jest rzeczywiście symetryczna.
\end{przyklad}

\section{Dlaczego średnie MCMC zbieżne są w szczególnych warunkach?}
\begin{twierdzenie}[Jednoznaczność i zbieżność z warunku minoracji]
\label{thm:mcmc-minorization}
Niech $P$ będzie jądrem Markowa na mierzalnej przestrzeni
$(\Theta,\mathcal A)$. Załóżmy, że istnieją
$m\ge1$, $\varepsilon\in(0,1]$ i prawdopodobieństwo $\nu$
takie, że dla każdego $x\in\Theta$ i $A\in\mathcal A$
$P^m(x,A)\ge\varepsilon\nu(A)$.
Wtedy $P$ ma dokładnie jedną miarę niezmienniczą $\pi$
i dla każdej początkowej miary $\mu$,
\[
 \norm{\mu P^{km}-\pi}_{\rm TV}
 \le(1-\varepsilon)^k\norm{\mu-\pi}_{\rm TV},
 \qquad k\ge0.
\]
W szczególności całkowita wariacja rozkładu łańcucha
dąży do zera również dla ogólnych $N\to\infty$.
\end{twierdzenie}
\begin{proof}
Jeżeli $0<\varepsilon<1$, zapisujemy
$P^m(x,\cdot)=\varepsilon\nu(\cdot)
 +(1-\varepsilon)R(x,\cdot)$, gdzie $R$ jest jądrem
Markowa. Dla dwóch miar probabilistycznych $\mu_1,\mu_2$
mamy
$\mu_1P^m-\mu_2P^m=(1-\varepsilon)(\mu_1R-\mu_2R)$.
Ponieważ jądro Markowa nie zwiększa odległości
w całkowitej wariacji, otrzymujemy kontrakcję o czynnik
$1-\varepsilon$. Przestrzeń miar probabilistycznych
z odległością całkowitej wariacji jest zupełna,
więc twierdzenie Banacha o punkcie stałym daje
jedyny stały punkt $\pi$ operatora $\mu\mapsto\mu P^m$.
Ponieważ $\pi P$ także jest jego punktem stałym,
$\pi P=\pi$. Każda miara niezmiennicza dla $P$
jest też niezmiennicza dla $P^m$, więc musi być $\pi$.
Iterowanie kontrakcji daje oszacowanie.
Dla $N=km+r$, $0\le r<m$, korzystamy z
$\pi P^r=\pi$ i nieekspansywności jądra $P^r$.
Przy $\varepsilon=1$ mamy $P^m(x,\cdot)=\nu$
dla wszystkich $x$ i tezy są natychmiastowe.
\end{proof}

\begin{uwaga}[Warunek jest mocny, ale przejrzysty]
Dla wielu ciągłych przestrzeni stanów, zwłaszcza
nieograniczonych, globalna minoracja nie zachodzi.
Używa się wtedy warunków Harrisa, funkcji Lyapunova
i małych zbiorów. Nie zakładamy, że każda propozycja
Metropolisa--Hastingsa daje zbieżność geometryczną.
Co więcej, samo trafienie w pobliże jednego maksimum
posterioru nie gwarantuje odwiedzenia innych modów.
\end{uwaga}

\begin{twierdzenie}[Prawo wielkich liczb dla skończonego łańcucha]
\label{thm:mcmc-finite-slln}
Niech $\Theta=\{1,\ldots,K\}$ i $P$ będzie macierzą
przejścia nieprzywiedlnego łańcucha Markowa na $\Theta$.
Niech $\pi$ oznacza jego jedyny rozkład stacjonarny.
Dla dowolnej funkcji $f\colon\Theta\to\R$ oraz
\emph{dowolnego} rozkładu początkowego,
$N^{-1}\sum_{k=0}^{N-1}f(X_k)\to\sum_{x=1}^Kf(x)\pi(x)$
prawie na pewno. Nieokresowość nie jest potrzebna
do tego prawa wielkich liczb.
\end{twierdzenie}
\begin{proof}
Niech $g=f-\pi(f)$, gdzie
$\pi(f)=\sum_xf(x)\pi(x)$. Dla skończonego,
nieprzywiedlnego łańcucha równanie Poissona
$(I-P)h=g$ ma rozwiązanie $h$, ponieważ przestrzeń
wektorów o średniej $\pi$ równej zero jest obrazem
$I-P$. Istotnie, wektor w jądrze $I-P$ jest stały:
zasada maksimum i nieprzywiedlność wymuszają
równość współrzędnych, a stąd rząd $I-P$ wynosi $K-1$.
Zatem $g(X_k)=h(X_k)-Ph(X_k)$ i
\[
 \sum_{k=0}^{N-1}g(X_k)
 =h(X_0)-h(X_N)+\sum_{k=1}^{N}
   \big(h(X_k)-Ph(X_{k-1})\big).
\]
Suma po prawej stronie jest martyngałem z przyrostami
o zerowej warunkowej średniej, a wszystkie przyrosty
są ograniczone stałą $2\norm h_\infty$.
Z nierówności Azumy--Hoeffdinga dla takich martyngałów
wynika, że dla każdego $\delta>0$ prawdopodobieństwo
przekroczenia przez wartość bezwzględną martyngału
progu $\delta N$ jest ograniczone przez
$2\exp[-c_\delta N]$ dla pewnego $c_\delta>0$.
Pierwszy lemat Borela--Cantellego daje zatem
zbieżność martyngału podzielonego przez $N$ do zera
prawie na pewno. Wyraz brzegowy jest ograniczony
przez $2\norm h_\infty/N\to0$. Sumując po
przeliczalnym zbiorze dodatnich wymiernych $\delta$
otrzymujemy tezę.
\end{proof}

\begin{lemat}[CLT dla ograniczonych różnic martyngałowych]
\label{lem:bounded-martingale-clt}
Niech $(D_k,\mathcal F_k)_{k\ge1}$ będzie ciągiem różnic
martyngałowych na jednej przestrzeni probabilistycznej,
$\E[D_k\mid\mathcal F_{k-1}]=0$, oraz niech
$|D_k|\le C$ prawie na pewno dla pewnego $C<\infty$
i wszystkich $k$. Oznaczmy
$V_k=\E[D_k^2\mid\mathcal F_{k-1}]$.
Jeżeli $N^{-1}\sum_{k=1}^NV_k\to v$ według
prawdopodobieństwa dla pewnej stałej $v\ge0$, to
$N^{-1/2}\sum_{k=1}^ND_k\Rightarrow\Normal(0,v)$.
Dla $v=0$ przez $\Normal(0,0)$ rozumiemy masę Diraca w zerze.
\end{lemat}
\begin{proof}
Ustalmy $t\in\R$ i napiszmy
$a_{N,k}=\E[\exp(itD_k/\sqrt N)\mid\mathcal F_{k-1}]$.
Z rozwinięcia wykładnika do drugiego rzędu, ograniczoności
$D_k$ i zerowej warunkowej średniej wynika jednostajnie
w $k\le N$ i niemal wszystkich realizacjach
\[
 a_{N,k}=1-\frac{t^2V_k}{2N}+O(N^{-3/2}),
 \quad
 \log a_{N,k}=-\frac{t^2V_k}{2N}+O(N^{-3/2}).
\]
Dla dużego $N$ wszystkie $a_{N,k}$ są niezerowe i ich
moduły są co najmniej $1-C_t/N$ dla pewnej stałej $C_t$.
Dla $S_k=\sum_{j=1}^kD_j$ proces
\[
 M_{N,k}=\frac{\exp(itS_k/\sqrt N)}
 {\prod_{j=1}^k a_{N,j}},\qquad 0\le k\le N,
\]
jest martyngałem zespolonym względem $\mathcal F_k$,
gdyż mianownik po dodaniu czynnika $a_{N,k}$ jest
$\mathcal F_{k-1}$-mierzalny. Mamy $\E M_{N,N}=1$,
a ponadto $|M_{N,N}|\le(1-C_t/N)^{-N}\le e^{2C_t}$
dla dostatecznie dużych $N$.
Zsumowanie rozwinięć logarytmu daje
$\prod_{k=1}^N a_{N,k}=
\exp[-t^2(2N)^{-1}\sum_{k=1}^NV_k+O(N^{-1/2})]$.
Ponieważ $N^{-1}\sum V_k\to v$ według prawdopodobieństwa,
ten iloczyn dąży do $\exp(-t^2v/2)$ według
prawdopodobieństwa i jest jednostajnie ograniczony.
Korzystając z ograniczoności $M_{N,N}$ otrzymujemy
\[
 \E\exp(itS_N/\sqrt N)
 =\E\left[M_{N,N}\prod_{k=1}^Na_{N,k}\right]
 \longrightarrow e^{-t^2v/2}\E M_{N,N}
 =e^{-t^2v/2}.
\]
Twierdzenie Lévy'ego o ciągłości funkcji charakterystycznych
kończy dowód.
\end{proof}

\section{Autokorelacja, długookresowa wariancja i ESS}
\begin{stwierdzenie}[Wariancja średniej ze stacjonarnego łańcucha]
\label{prop:mcmc-autocorrelation}
Niech $(X_k)_{k\ge0}$ będzie ściśle stacjonarnym
łańcuchem Markowa o rozkładzie $\pi$ i niech
$Y_k=f(X_k)$ ma skończoną, dodatnią wariancję
$\gamma_0=\Var_\pi(Y_0)$. Oznaczmy
$\gamma_h=\Cov_\pi(Y_0,Y_h)$ oraz
$\rho_h=\gamma_h/\gamma_0$. Wtedy dla każdego $N\ge1$
\[
 \Var\left(\frac1N\sum_{k=0}^{N-1}Y_k\right)
 =\frac{\gamma_0}{N}\left(1+2\sum_{h=1}^{N-1}
 (1-h/N)\rho_h\right).
\]
Jeżeli $\sum_{h\ge1}|\gamma_h|<\infty$, to
$N\Var(N^{-1}\sum_{k=0}^{N-1}Y_k)\to
\sigma_\infty^2=\gamma_0+2\sum_{h\ge1}\gamma_h\ge0$.
Gdy $\sigma_\infty^2>0$, definiujemy
$\tau_{\rm int}=\sigma_\infty^2/\gamma_0$ i
$N_{\rm eff}=N/\tau_{\rm int}$ jako asymptotyczną
\emph{efektywną liczebność próby dla funkcji $f$}.
\end{stwierdzenie}
\begin{proof}
Wariancja sumy wynosi
$\sum_{i,j=0}^{N-1}\Cov(Y_i,Y_j)$.
Ze stacjonarności, dla każdego opóźnienia $h\ge1$
występuje $2(N-h)$ par o kowariancji $\gamma_h$
oraz $N$ par z kowariancją $\gamma_0$.
Podzielenie przez $N^2$ daje pierwszy wzór.
Przy bezwzględnej sumowalności można zastosować
twierdzenie o zbieżności zdominowanej dla szeregu,
ponieważ $|1-h/N|\le1$ dla $h<N$.
Nieujemność granicy wynika z nieujemności wariancji.
\end{proof}

\begin{twierdzenie}[Skończony odwracalny łańcuch: widmo i CLT]
\label{thm:mcmc-reversible-clt}
Niech $P$ będzie nieprzywiedlną i nieokresową
macierzą przejścia na skończonej przestrzeni $\Theta$,
odwracalną względem stacjonarnego $\pi$.
Dla $f\colon\Theta\to\R$ o $\pi$-średniej $\mu$ niech
$g=f-\mu$. Niech $(\lambda_j,e_j)_{j=2}^K$
stanowi układ ortonormalnych wektorów własnych
operatora $P$ w $L^2(\pi)$ prostopadłych do stałych,
tak że $g=\sum_{j=2}^Kc_je_j$.
Wtedy $-1<\lambda_j<1$ i
\[
 \sigma_\infty^2
 =\sum_{j=2}^K c_j^2\frac{1+\lambda_j}{1-\lambda_j}.
\]
Jeżeli $g\not\equiv0$, to $\sigma_\infty^2>0$.
Przy \emph{dowolnym} rozkładzie początkowym zachodzi
\[
 \sqrt N\left(\frac1N\sum_{k=0}^{N-1}f(X_k)-\mu\right)
 \Rightarrow\Normal(0,\sigma_\infty^2).
\]
\end{twierdzenie}
\begin{proof}
Szczegółowa równowaga oznacza, że $P$ jest samosprzężony
w $L^2(\pi)$. Jest operatorem kontrakcji, więc jego
wartości własne leżą w $[-1,1]$. Nieprzywiedlność i
nieokresowość w skończonym stanie dają prostą wartość
własną $1$ i wykluczają $-1$.
Dla startu stacjonarnego
$\gamma_h=\langle g,P^hg\rangle_\pi=
\sum_{j=2}^Kc_j^2\lambda_j^h$.
Sumując zbieżny szereg geometryczny otrzymujemy
$\gamma_0+2\sum_{h\ge1}\gamma_h
=\sum_jc_j^2(1+\lambda_j)/(1-\lambda_j)$.
Aby uzasadnić CLT, rozwiązujemy $(I-P)h=g$ tak jak
w dowodzie twierdzenia~\ref{thm:mcmc-finite-slln}.
Suma $g(X_k)$ różni się o ograniczony wyraz brzegowy
od sumy różnic martyngałowych
$D_k=h(X_k)-Ph(X_{k-1})$. Ich warunkowe wariancje
wynoszą $V(X_{k-1})=Ph^2(X_{k-1})-(Ph(X_{k-1}))^2$.
Z udowodnionego prawa wielkich liczb
$N^{-1}\sum_{k=1}^NV(X_{k-1})\to\pi(V)$
prawie na pewno. Przyrosty $D_k$ są ograniczone,
więc spełniają warunek Lindeberga. Stosujemy lemat~\ref{lem:bounded-martingale-clt},
którego założenia zostały tu sprawdzone.
Otrzymujemy normalną granicę z wariancją $\pi(V)$. Ta wariancja
równa się obliczonej wcześniej $\sigma_\infty^2$,
gdyż oba wzory są granicą $N^{-1}\Var_\pi(
\sum_{k=0}^{N-1}g(X_k))$; dla ograniczonych przyrostów
różnice od wyrazu brzegowego są $o(N)$ w $L^2$.
Wyjściowy rozkład nie wpływa na CLT: po skończonej
liczbie kroków jego wpływ na znormalizowaną sumę
zanika, co można formalnie uzasadnić sprzężeniem
ze stacjonarną kopią łańcucha (zbieżność geometryczna
w skończonej nieprzywiedlnej nieokresowej przestrzeni).
\end{proof}

\begin{uwaga}[Czego ESS nie gwarantuje]
Dla silnie dodatniej autokorelacji $\tau_{\rm int}\gg1$,
więc $N_{\rm eff}\ll N$. Przy ujemnych korelacjach
może być też $\tau_{\rm int}<1$, a $N_{\rm eff}>N$;
nie jest to liczba faktycznie niezależnych obserwacji.
Efektywna liczebność zależy od funkcji $f$ i od
stacjonarnego rozkładu. Oszacowanie ESS z krótkiego,
źle mieszającego się łańcucha może być bardzo mylące.
\end{uwaga}

\section{Gibbs sampling i modele hierarchiczne}
\begin{twierdzenie}[Kroki Gibbsa zachowują posterior]
\label{thm:gibbs-invariant}
Niech $\pi(x,y)$ będzie wspólną, właściwą gęstością
na $\R^r\times\R^s$ względem $\lambda_{r+s}$.
Załóżmy, że warunkowe gęstości $\pi(x\mid y)$ oraz
$\pi(y\mid x)$ są określone dla prawie wszystkich
wartości odpowiednich zmiennych. Jądro aktualizujące
$x$ przez losowanie z $\pi(\cdot\mid y)$ przy
ustalonym $y$ oraz jądro aktualizujące $y$ przez
losowanie z $\pi(\cdot\mid x)$ przy ustalonym $x$
mają $\pi$ jako miarę niezmienniczą. Ich złożenie,
\emph{systematyczny skan Gibbsa}, również zachowuje $\pi$.
Nie musi jednak samo być odwracalne względem $\pi$.
\end{twierdzenie}
\begin{proof}
Dla zbiorów mierzalnych $A\subset\R^r$, $B\subset\R^s$
po aktualizacji pierwszej współrzędnej otrzymujemy
\[
 \int_{\R^s}\pi_Y(y)
 \left(\int_A\pi(x\mid y)\,\dd\lambda_r(x)\right)
 \ind{B}(y)\,\dd\lambda_s(y)
 =\int_{A\times B}\pi(x,y)\,\dd\lambda_{r+s}(x,y).
\]
Jest to dokładnie $\pi(A\times B)$. Zbiory prostokątne
generują borelowskie $\sigma$-ciało, więc pierwsze
jądro zachowuje $\pi$; identycznie rozumujemy dla drugiego.
Z równości $\pi P_x=\pi$ i $\pi P_y=\pi$
wynika $\pi(P_xP_y)=\pi$.
\end{proof}

\begin{przyklad}[Gibbs dla skorelowanej pary normalnej]
Niech $(X,Y)$ ma normalny rozkład dwuwymiarowy o średniej
$(0,0)$, wariancjach $1$ i korelacji $\rho\in(-1,1)$.
Wtedy
$X\mid Y=y\sim\Normal(\rho y,1-\rho^2)$ oraz
$Y\mid X=x\sim\Normal(\rho x,1-\rho^2)$.
Algorytm naprzemiennie losuje z tych dwóch rozkładów.
Gdy $|\rho|$ jest bliskie jedności, kolejne iteracje
są mocno powiązane, mimo że \emph{każdy krok warunkowy}
ma dokładny rozkład. Może być wtedy lepiej wykonać
aktualizację blokową lub zmienić parametryzację.
\end{przyklad}

\section{Hamiltonian Monte Carlo: dlaczego gradient pomaga?}
\label{sec:hmc-intro}
Niech posterior ma dodatnią, różniczkowalną gęstość
$\pi(q)\propto\exp[-U(q)]$ na $\R^d$, a $M$
będzie dodatnio określoną macierzą masy. Wprowadzamy
pomocniczy pęd $p\sim\Normal_d(0,M)$ oraz
\emph{hamiltonian} $H(q,p)=U(q)+\frac12p^\top M^{-1}p$.
Wspólny rozkład docelowy na przestrzeni fazowej
jest proporcjonalny do $\exp[-H(q,p)]$.
Ideą jest wykonywanie długich, skorelowanych ruchów
wzdłuż przybliżonych trajektorii hamiltonowskich,
zamiast małych, losowych kroków bez informacji o gradiencie.

\begin{stwierdzenie}[Zachowanie rozkładu przez krok HMC]
\label{prop:hmc-step}
Załóżmy, że $U\in C^1(\R^d)$, $M$ jest dodatnio
określona i $\int_{\R^d}e^{-U(q)}\,\dd\lambda_d(q)<\infty$.
Niech $\Phi$ będzie złożeniem skończonej liczby
kroków algorytmu leapfrog o stałym kroku $\epsilon>0$:
\[
 \begin{aligned}
 p_{1/2}&=p-\tfrac\epsilon2\nabla U(q),\\
 q'&=q+\epsilon M^{-1}p_{1/2},\\
 p'&=p_{1/2}-\tfrac\epsilon2\nabla U(q').
 \end{aligned}
\]
Po złożeniu z odwróceniem pędu $R(q,p)=(q,-p)$
propozycja $T=R\circ\Phi$ jest odwracalną inwolucją
zachowującą miarę Lebesgue'a na $\R^{2d}$.
Akceptowanie jej z prawdopodobieństwem
$\min\{1,\exp[H(q,p)-H(T(q,p))]\}$ zachowuje
wspólny rozkład o gęstości proporcjonalnej do $e^{-H}$.
Odświeżenie pędu z $\Normal_d(0,M)$ również zachowuje
ten rozkład. Łańcuch położeń $q$ ma więc $\pi$
jako rozkład niezmienniczy.
\end{stwierdzenie}
\begin{proof}
Każda podoperacja leapfrog jest ścinaniem:
$(q,p)\mapsto(q,p-\frac\epsilon2\nabla U(q))$
lub $(q,p)\mapsto(q+\epsilon M^{-1}p,p)$.
Dla ustalonej drugiej współrzędnej pierwsze ścinanie
przesuwa $p$, a drugie dla ustalonego $p$ przesuwa $q$;
obie transformacje zachowują miarę Lebesgue'a
(bez potrzeby zakładania istnienia drugich pochodnych).
Każde ma odwrotność przez zmianę $\epsilon$ na $-\epsilon$.
Symetria kompozycji daje $R\Phi R=\Phi^{-1}$,
stąd $(R\Phi)^2=I$. Również $R$ zachowuje objętość.
Dla odwracalnej, zachowującej objętość inwolucji $T$
przepływ masy między $z$ i $T(z)$ po odrzuceniach
jest symetryczny:
$e^{-H(z)}\min(1,e^{H(z)-H(Tz)})
=\min(e^{-H(z)},e^{-H(Tz)})$.
Jest to szczegółowa równowaga na przestrzeni fazowej.
Odświeżenie pędu zachowuje produkt rozkładu $q$ i
normalnego rozkładu $p$, zatem złożenie zachowuje
wspólny rozkład. Po marginalizacji otrzymujemy $\pi$.
\end{proof}

\begin{uwaga}[Co wymaga kontroli przy HMC]
Samo zachowanie rozkładu nie dowodzi ergodyczności.
Błędny gradient, niestabilne kroki, zbyt duże $\epsilon$,
wielomodalność lub parametryzacja z wąskimi obszarami
mogą prowadzić do problemów. Należy kontrolować
akceptację, rozbieżne trajektorie numeryczne, efektywne
liczebności próby i zgodność wielu łańcuchów.
Właściwa implementacja NUTS wymaga osobnego uzasadnienia
reguły zatrzymywania trajektorii; nie utożsamiamy jej
z opisanym tu algorytmem o stałej liczbie kroków.
\end{uwaga}

\section{Diagnostyka: co można sprawdzić, a czego nie można udowodnić wykresem?}
\begin{definicja}[Dwie proste diagnostyki wielu łańcuchów]
Załóżmy, że mamy $m\ge2$ łańcuchów, każdy długości
$N\ge2$, wartości $z_{ji}$ badanego funkcjonału,
$\overline z_j=N^{-1}\sum_{i=1}^Nz_{ji}$ oraz
$\overline z=m^{-1}\sum_{j=1}^m\overline z_j$.
Definiujemy
\[
 W=\frac1m\sum_{j=1}^m
 \frac1{N-1}\sum_{i=1}^N(z_{ji}-\overline z_j)^2,
 \qquad
 B=\frac N{m-1}\sum_{j=1}^m(\overline z_j-\overline z)^2.
\]
Jeżeli $W>0$, to klasyczna diagnostyka Gelmana--Rubina
ma postać
\[
 \widehat R=\sqrt{\frac{(N-1)W+B}{NW}}.
\]
W praktyce warto stosować podział łańcuchów na połowy
(\emph{split-$\widehat R$}) oraz wersje rangowe
i diagnostykę ogonów rozkładu.
\end{definicja}

\begin{uwaga}[Wygrzewanie i przerzedzanie łańcucha]
Odrzucenie początkowych iteracji (\emph{warm-up}, także \emph{burn-in})
może zmniejszać wpływ startu, ale arbitralna liczba odrzuconych kroków
nie gwarantuje osiągnięcia stacjonarności. Przerzedzanie (\emph{thinning}),
czyli zapisywanie co $k$-tej iteracji, na ogół nie poprawia
precyzji estymacji przy ustalonym koszcie wygenerowania pełnego łańcucha;
może natomiast ograniczać zużycie pamięci. W raporcie należy podać
liczbę iteracji przed i po warm-up, liczbę łańcuchów i sposób diagnostyki.
\end{uwaga}

\begin{przyklad}[Pozorna zgodność łańcuchów i nieodwiedzony mod]
Niech docelowy rozkład na $\R$ będzie mieszaniną
$\tfrac12\Normal(-a,1)+\tfrac12\Normal(a,1)$ dla dużego $a$.
Jeżeli wszystkie łańcuchy rozpoczynamy w pobliżu $-a$ i używamy
małych kroków random-walk Metropolisa, mogą przez cały dostępny
czas odwiedzać wyłącznie lewy mod. Ich średnie i wariancje
międzyłańcuchowe będą podobne, zatem $\widehat R$ może być
bliska jedności, mimo że estymata $\E_\pi X=0$ będzie bliska $-a$.
To nie jest kontrprzykład do poprawności algorytmu Metropolisa,
lecz do \emph{wnioskowania o zbieżności jedynie z diagnostyki}.
Należy stosować różne punkty startowe, w tym odległe mody,
i sprawdzać stabilność oszacowań przy zmianach parametryzacji
oraz propozycji.
\end{przyklad}

\begin{uwaga}[Nie istnieje magiczny próg zapewniający zbieżność]
Wartość $\widehat R$ bliska jeden nie dowodzi,
że łańcuchy odwiedziły wszystkie istotne mody. Może
się zdarzyć, że wszystkie łańcuchy utknęły w tym samym
obszarze. Same wykresy śladu i histogramy również nie
stanowią dowodu osiągnięcia rozkładu stacjonarnego.
Bezpieczna procedura obejmuje różne punkty startowe,
kontrolę $\widehat R$, ESS dla średnich i ogonów,
błąd Monte Carlo, analizę posterior predictive oraz
badanie wrażliwości na parametry algorytmu.
\end{uwaga}

\begin{center}\small
\begin{tabular}{>{\raggedright\arraybackslash}p{.25\textwidth}>{\raggedright\arraybackslash}p{.30\textwidth}>{\raggedright\arraybackslash}p{.29\textwidth}}
\toprule
\textbf{Sytuacja} & \textbf{Algorytm} & \textbf{Na co uważać}\\
\midrule
Dokładne warunkowe rozkłady łatwe do losowania & Gibbs & Zależność współrzędnych, wolne mieszanie\\
Posterior jednowymiarowy, gęstość do stałej & Random-walk MH & Skala propozycji, autokorelacja\\
Gładki posterior w większym wymiarze & HMC/NUTS & Gradient, rozbieżności, parametryzacja\\
Trudny multimodalny posterior & Metody tempering / specjalne propozycje & Pozorne zbieganie wszystkich łańcuchów\\
Proste losowanie niezależne dostępne & Zwykłe MC, nie MCMC & Po co wprowadzać zależność?\\
\bottomrule
\end{tabular}
\end{center}

\begin{center}
\begin{tikzpicture}[>=Latex, every node/.style={font=\small},node distance=16mm]
\node[draw,rounded corners,fill=blue!5,minimum width=22mm] (a1) {Losowanie 1};
\node[draw,rounded corners,fill=blue!5,minimum width=22mm,right=of a1] (a2) {Losowanie 2};
\node[draw,rounded corners,fill=blue!5,minimum width=22mm,right=of a2] (a3) {Losowanie 3};
\node[below=12mm of a1,draw,rounded corners,fill=green!6] (b1) {$X_0$};
\node[below=12mm of a2,draw,rounded corners,fill=green!6] (b2) {$X_1$};
\node[below=12mm of a3,draw,rounded corners,fill=green!6] (b3) {$X_2$};
\draw[->,thick] (b1)--node[above] {$P$} (b2);
\draw[->,thick] (b2)--node[above] {$P$} (b3);
\node[below=7mm of b2,font=\footnotesize,text width=100mm,align=center]
 {Górny wiersz: niezależne losowania z rozkładu docelowego.\\
 Dolny wiersz: łańcuch MCMC, zależność przez jądro $P$.};
\end{tikzpicture}

{\small\emph{Rysunek 24.1.} Zwykłe Monte Carlo i MCMC wymagają odmiennych wzorów na błąd standardowy.}
\end{center}

\section{Laboratorium: posterior beta, cztery łańcuchy i diagnostyka}
\label{sec:mcmc-laboratorium}
Na danych o $7$ wadach ze $100$ prób przy priorze
$\operatorname{Beta}(1,1)$ posterior ma rozkład
$\operatorname{Beta}(8,94)$. Ponieważ umiemy z niego
losować \emph{bezpośrednio}, w rzeczywistym zastosowaniu
wybralibyśmy samodzielne losowania Monte Carlo.
Celowo użyjemy jednak tego przykładu do sprawdzenia
implementacji Metropolisa--Hastingsa: wiemy, jaka
jest dokładna średnia posterioru, i możemy porównać
wynik łańcuchów z wartością $8/102$.
Współrzędna $z=\log[p/(1-p)]$ obejmuje całą $\R$.
Ważne: po transformacji gęstość względem $\dd\lambda_1(z)$
zawiera \emph{jakobian} $p(1-p)$.
Stąd dla $p=(1+e^{-z})^{-1}$ log-gęstość docelowa
współrzędnej $z$ jest (z dokładnością do stałej)
$a\log p+b\log(1-p)$, a nie
$(a-1)\log p+(b-1)\log(1-p)$.

\begin{lstlisting}[language=Python,caption={Kompletny program: cztery łańcuchy Metropolisa--Hastingsa dla posterioru beta, split-R-hat i przybliżony ESS (Python, NumPy i SciPy).}]
import numpy as np
from scipy.special import expit

A, B = 8.0, 94.0
N, BURN, CHAINS = 24_000, 4_000, 4
STEP = 0.65
rng = np.random.default_rng(2026)

def log_target(z):
    # Posterior in logit coordinate: beta density times Jacobian.
    log_p = -np.logaddexp(0.0, -z)
    log_1mp = -np.logaddexp(0.0, z)
    return A * log_p + B * log_1mp

def run_chain(z0):
    z, accepted = float(z0), 0
    samples = np.empty(N)
    current = log_target(z)
    for k in range(N):
        proposed = z + rng.normal(0.0, STEP)
        newlog = log_target(proposed)
        if np.log(rng.random()) < min(0.0, newlog - current):
            z, current = proposed, newlog
            accepted += 1
        samples[k] = expit(z)
    return samples[BURN:], accepted / N

starts = np.array([-5.0, -3.0, -1.0, 1.0])
results = [run_chain(z0) for z0 in starts]
chains = np.stack([r[0] for r in results])
print("acceptances:", [round(r[1], 3) for r in results])
print("posterior mean estimate:", chains.mean())
print("exact posterior mean:", A / (A + B))

# Split R-hat (original moment-based version).
L = chains.shape[1]
half = L // 2
split = np.concatenate([chains[:, :half], chains[:, -half:]], axis=0)
m = split.shape[0]
n = split.shape[1]
means = split.mean(axis=1)
W = split.var(axis=1, ddof=1).mean()
B_between = n * means.var(ddof=1)
rhat = np.sqrt(((n - 1) * W / n + B_between / n) / W)
print("split R-hat:", rhat)

# Descriptive autocorrelation-based ESS (not a convergence proof).
def acf_fft(x):
    x = np.asarray(x) - np.mean(x)
    n = len(x)
    cov = np.fft.irfft(np.abs(np.fft.rfft(x, n=2*n))**2,
                       n=2*n)[:n]
    return cov / cov[0]

acf = np.array([acf_fft(c) for c in chains])
rho = acf.mean(axis=0)
tau = 1.0
# Truncate at the first nonpositive adjacent-pair sum.
for k in range(1, min(1000, len(rho) - 1), 2):
    pair = rho[k] + rho[k + 1]
    if pair <= 0:
        break
    tau += 2 * pair
print("approximate ESS:", chains.size / tau)
print("approximate MC SE:", np.sqrt(chains.var(ddof=1)
                                     * tau / chains.size))
\end{lstlisting}

\begin{lstlisting}[language=Julia,caption={Kompletny program: random-walk Metropolis dla posterioru beta we współrzędnej logit (Julia, Random i Statistics).}]
using Random, Statistics
Random.seed!(2026)
a, b = 8.0, 94.0
N, burn, step = 30_000, 5_000, 0.65

sigmoid(z) = z >= 0 ? inv(1 + exp(-z)) : exp(z) / (1 + exp(z))
softplus(z) = max(z, 0.0) + log1p(exp(-abs(z)))
logtarget(z) = -a * softplus(-z) - b * softplus(z)

z = log(a / b)
oldlog = logtarget(z)
samples = zeros(Float64, N)
accept = 0
for k in 1:N
    prop = z + step * randn()
    newlog = logtarget(prop)
    if log(rand()) < min(0.0, newlog - oldlog)
        z, oldlog = prop, newlog
        accept += 1
    end
    samples[k] = sigmoid(z)
end
println("acceptance: ", accept / N)
println("MCMC posterior mean: ", mean(samples[(burn+1):end]))
println("exact posterior mean: ", a / (a+b))
\end{lstlisting}

\begin{zadanie}[Zależność zmienia wariancję średniej]
Załóżmy, że w stacjonarnym łańcuchu $\Var(Y_0)=4$ oraz
$\rho_h=\rho^h$ dla pewnego $|\rho|<1$.
Wyznaczyć $\sigma_\infty^2$ i asymptotyczny ESS.
\par\emph{Rozwiązanie.} Z sumy szeregu geometrycznego
$\tau_{\rm int}=1+2\rho/(1-\rho)=(1+\rho)/(1-\rho)$,
$\sigma_\infty^2=4(1+\rho)/(1-\rho)$, a
$N_{\rm eff}=N(1-\rho)/(1+\rho)$.
Przy $\rho=0{,}9$ efektywna liczebność jest w przybliżeniu
$N/19$, chociaż samych iteracji jest $N$.
Dla ujemnej korelacji ESS może przewyższać $N$.
\end{zadanie}

